%=======================================================================
% Requires: booktabs, multirow, amsmath, amssymb, siunitx, todonotes
% Positioning against prior datasets is in Section~\ref{sec:related}.
%=======================================================================
\section{Dataset Design and Characterization}
\label{sec:collection}

\subsection{Collection Protocol}

\subsubsection{Sites and Roles}
\label{sec:environments}

Data collection was conducted on the Guangfu Campus of National Yang Ming Chiao
Tung University (NYCU), at indoor and outdoor sites. Site selection follows
three factors, each tied to a quantity the release reports per frame rather than
to a qualitative descriptor: \emph{viewpoint}, whether the site admits vantages
displaced from one another---an elevated infrastructure position, and mobile
nodes able to separate in the ground plane---characterised by the range of
aspect angles $\beta$ realisable between an ego and a helper;
\emph{appearance}, the illumination, texture, and clutter that set the range
beyond which $\mathrm{visible}(i)$ ceases to imply $\mathrm{observed}(i)$; and
\emph{radio propagation}, whether structure intervenes between a node and the
access point, characterised by first-Fresnel-zone clearance. Each site is
surveyed and mapped once as a dense point cloud and every node pose expressed in
one anchored map frame, so that occlusion, co-visibility, and link geometry are
evaluated against a single geometric representation. Within a run, one mobile
agent is the \emph{ego}, viewpoint limited by design---routed so that structure
intervenes between it and the annotated objects, or carrying the shorter-range
depth sensor---and the other the \emph{partner}; ego assignment alternates
across runs, so that a reported condition is attributable to the geometry rather
than to a platform. Infrastructure nodes are never ego: their function is the
elevated, persistent vantage no mobile node sustains, and assigning one as ego
would reinstate the common-height symmetry this collection is designed to break.
The sequence characterised below was recorded at the MIRC site, an indoor
laboratory with one surveyed elevated camera--radar mast, a
\SI{5.75}{\giga\hertz} access point positioned behind furniture, and static
seating that constitutes the occluding structure.

\todo{Site-overview figure: map cloud, surveyed infrastructure pose, AP
position, and both mobile trajectories, one panel per site.}

\subsubsection{Collaboration Geometry}
\label{sec:paradigms}

A paradigm specifies a trajectory geometry together with the condition that
geometry is intended to induce. The vocabulary is that of paradigm-driven
multi-robot collection; what is added is that each paradigm is a predicate over
per-frame quantities, computed from the map, the annotations, and the node
poses, and released with the recordings rather than asserted of a run.

Consider one annotated object at one frame. Let $e$ denote the ego and $j$ range
over the remaining nodes. For node $i$, let $\omega_i \in [0,1]$ be the fraction
of rays cast from $i$'s sensor origin to the object surface that are blocked by
mapped structure, $r_i$ the range to the object, and $n_i$ the evidence returns
falling within the annotated box---LiDAR points, depth points, or projected
pixel area, according to the node's primary modality. With $\tau = 0.5$,
%
\begin{align}
\mathrm{visible}(i) &\equiv \text{in FoV} \;\wedge\; \omega_i < \tau,
\label{eq:visible}\\
\mathrm{observed}(i) &\equiv \mathrm{visible}(i) \;\wedge\; n_i \ge n_{\min,i}
\;\wedge\; r_i \le r_{\max,i}.
\label{eq:observed}
\end{align}
%
Visibility is geometric---an unobstructed line exists---while observation is
evidential: the sensor returned support sufficient for a detection, within the
range at which that modality still returns anything. Both have single-agent
precedent (Section~\ref{sec:related}); here they are computed for every node and
kept distinct, so that an object the ego could not see remains separable from
one it could see but for which it returned too little evidence, a separation
the rescue columns of Section~\ref{sec:descriptor} depend on. A node $j \ne e$ is a
\emph{helper} if $\mathrm{visible}(j)$, and the best helper $j^\star$ is that of
lowest $\omega$; the aspect angle $\beta$ subtended at the object between the
rays to $e$ and to $j^\star$ distinguishes a redundant viewpoint near $0^\circ$
from a complementary one beyond $90^\circ$. Each frame further carries the
co-visibility fraction $C$, objects visible to at least two nodes over those
visible to at least one, and, per link, the minimum first-Fresnel-zone clearance
against the same occupancy grid normalised by the Fresnel radius, with
$\mathrm{clear}(a,b) < 0.6$ the obstructed-link criterion.
Table~\ref{tab:paradigms} composes these into the paradigms, evaluated in order
so that a condition combining occlusion with a degraded link is reported as such
rather than absorbed into the plain occlusion class; there $v_o$ is the object
speed, $v_{r,i}$ its radial velocity relative to node $i$, and
$v_{\min} = k\lambda / (2 N_c T_c)$ with $k \approx 1.5$ the Doppler resolution
floor of the radar.

\begin{table}[t]
\centering
\caption{Trajectory paradigms, evaluated in order; the first match is assigned.
$\tau = 0.5$; helper and $j^\star$ as defined in Section~\ref{sec:paradigms}.
Superscripts state what the present release supports: $^{\checkmark}$ computed
and released, $^{\circ}$ computed but not folded into the classifier,
$^{\times}$ undefined in the absence of moving targets.}
\label{tab:paradigms}
\footnotesize
\setlength{\tabcolsep}{4pt}
\begin{tabular}{@{}cll@{}}
\toprule
\# & Paradigm & Condition \\
\midrule
1 & Doppler diversity$^{\times}$      & $\lVert v_o\rVert > 0.3 \;\wedge\; \lVert v_{r,e}\rVert < v_{\min} \le \max_j \lVert v_{r,j}\rVert$ \\
2 & Occlusion $\times$ link$^{\circ}$ & $\neg\,\mathrm{visible}(e) \;\wedge\; \text{helper} \;\wedge\; \mathrm{clear}(e,j^\star) < 0.6$ \\
3 & Occlusion$^{\checkmark}$          & $\neg\,\mathrm{visible}(e) \;\wedge\; \text{helper}$ \\
4 & Handover$^{\checkmark}$           & $\lvert \{ i : \mathrm{visible}(i) \} \rvert = 1$ \\
5 & Link stress$^{\circ}$             & $\text{helper} \;\wedge\; \mathrm{clear}(e,j^\star) < 0.6$ \\
6 & Co-observation$^{\checkmark}$     & $\mathrm{visible}(e) \;\wedge\; \text{helper}$ \\
7 & Unstructured$^{\checkmark}$       & none of the above \\
\bottomrule
\end{tabular}
\end{table}

A label is assigned per frame from the measured metrics, not declared once for a
run: a sequence intended as co-observation carries occlusion labels wherever the
geometry produced them. Evaluation strata are therefore cut by $\omega$,
$\beta$, $r$, $C$, and clearance directly, never by the label. Each run also
records an \texttt{intended\_paradigm}, and comparing it against the measured
distribution shows how reliably a trajectory plan produced the condition it was
written for.

Table~\ref{tab:params} lists the pipeline parameters, frozen across sequences.
Three are workarounds rather than choices: the density threshold and
consecutive-hit test suppress speckle in the registered map, and the endpoint
clearances exist because the nodes and the access point are themselves mapped as
structure, so without them every link ray is blocked at its own endpoint. A
survey-grade map should permit their removal.

\begin{table}[t]
\centering
\caption{Pipeline parameters, frozen across sequences. Range caps are set from
the range at which median in-box evidence falls to its floor, then fixed.}
\label{tab:params}
\footnotesize
\begin{tabular}{@{}lll@{}}
\toprule
& Parameter & Value \\
\midrule
\multirow{5}{*}{Occlusion}
  & threshold $\tau$              & 0.5 \\
  & occupancy voxel               & \SI{10}{\centi\metre} \\
  & rays per object, node, frame  & 200 \\
  & outlier removal               & 20 neighbours, $2\sigma$ \\
  & occupied voxel / blocked ray  & $\ge 8$ points / 2 consecutive \\
\midrule
\multirow{3}{*}{Evidence}
  & floors                        & 1 LiDAR pt, 50 depth pts, 400\,px$^2$ \\
  & range caps                    & 20 / 4.5 / 5 / \SI{20}{\metre} \\
  &                               & (Ouster, ZED, RealSense, infra) \\
\midrule
\multirow{2}{*}{Link}
  & obstructed below              & $0.6\,r_1$ \\
  & carrier frequency             & \SI{5.75}{\giga\hertz}, read per frame \\
\midrule
\multirow{2}{*}{Workarounds}
  & sensor self-clearance         & \SI{0.6}{\metre} \\
  & link endpoint clearance       & \SI{0.7}{\metre} \\
\bottomrule
\end{tabular}
\end{table}

\subsection{Sequence Characterization}
\label{sec:descriptor}

Each sequence is summarised by one row per ego assignment in four blocks of
frozen columns (Table~\ref{tab:seq}). The \textit{design} block records what was
declared before the run: sequence, site, intended paradigm, the stressor applied
if any---low light, crowd, node outage, or induced packet loss, each a
controlled departure from a baseline run of the same trajectory---and which
node is ego. The remaining blocks record what was measured. The
\textit{geometry} block gives the median $C$, the median $\beta$ to
infrastructure, the proportion of observations at $\beta > 90^\circ$, and the
count of occlusion episodes, an episode being a contiguous interval of
$\neg\,\mathrm{visible}(e)$ on one object lasting at least \SI{0.5}{\second},
with the number covered throughout by a helper in parentheses. The
\textit{rescue} block gives the proportion of ego-in-FoV observations in which a
partner supplies what the ego does not: \textit{occlusion rescue}, the ego
blocked and a helper not; \textit{range rescue}, the ego holding an unobstructed
line but returning insufficient evidence,
$\mathrm{visible}(e) \wedge \neg\,\mathrm{observed}(e)$, while a helper
observes; and \textit{Doppler rescue}, the ego's radar blind to motion another
node's radar resolves. The final block gives the fraction of frames on which the
ego--AP link is obstructed or degraded, and the median number of annotated
objects in view. A quantity undefined for a sequence is left empty rather than
reported as zero.

\begin{table*}[t]
\centering
\caption{Per-sequence descriptor, shown for \texttt{coop2\_0828} under both ego
assignments. $^\ast$ staged ego. Median episode duration \SI{18.8}{\second};
elevation offset to the infrastructure node \SI{1.53}{\metre} (IQR
\SIrange{1.52}{1.54}{\metre}). Range rescue is reported under RGB-D evidence.
An em-dash denotes a quantity undefined for the sequence, not a zero.}
\label{tab:seq}
\small
\setlength{\tabcolsep}{4pt}
\begin{tabular}{@{}lllllccccccccc@{}}
\toprule
\multicolumn{5}{c}{Design} & \multicolumn{4}{c}{Geometry} &
\multicolumn{3}{c}{Rescue rate} & \multicolumn{2}{c}{Link, scene} \\
\cmidrule(r){1-5}\cmidrule(lr){6-9}\cmidrule(lr){10-12}\cmidrule(l){13-14}
Sequence & Site & Paradigm & Stressor & Ego &
$C$ & $\beta$ ($^\circ$) & $\beta > 90^\circ$ & Episodes &
Occl. & Range & Dopp. &
NLoS (\%) & Obj. \\
\midrule
\multirow{2}{*}{\texttt{coop2\_0828}} & \multirow{2}{*}{MIRC} &
\multirow{2}{*}{co-observation} & \multirow{2}{*}{none} &
\texttt{m1}$^\ast$ & 0.5 & 149 & 0.67 & 9 (6) & 0.10 & 0.29 & --- & 33 & 2 \\
 & & & & \texttt{m2} & 0.5 & 29 & 0.33 & 8 (5) & 0.02 & 0.39 & --- & 18 & 2 \\
\bottomrule
\end{tabular}
\end{table*}

The present release contains one sequence, \texttt{coop2\_0828}: \num{733}
synchronised frames over \SI{156}{\second} at the MIRC site, two mobile nodes
and one infrastructure node, two annotated static objects, reported under both
ego assignments. It was planned as co-observation and measures as
co-observation in 74\,\% and 77\,\% of ego-in-FoV observations respectively;
the remainder are occlusion, handover, and unstructured frames the plan did not
call for and the labels record regardless. Doppler diversity is left empty: the
predicate is undefined without moving targets, and the sequence contains none.

The row is read as an instruction for how the sequence can be used. The rescue
columns locate where collaboration is needed: range rescue at 0.29 and 0.39
against occlusion rescue at 0.10 and 0.02 says that a fusion method evaluated
here is tested mainly on frames where the ego sees the object but returns too
little evidence, not on frames where it is blocked, so an occlusion-recovery
claim should not be made from this sequence. The $\beta$ columns say which ego
assignment tests what: with \texttt{mobile\_1} as ego the infrastructure is a
complementary viewpoint on two frames in three, with \texttt{mobile\_2} on one
in three, so the same recording supports a wide-baseline and a redundant-view
evaluation depending on which node is scored. The NLoS column gives the
fraction of frames---33\,\% and 18\,\%---on which a link-aware method can be
evaluated under an obstructed channel, and $C = 0.5$ bounds what any late
fusion could recover per frame. Episodes, with their durations, are the unit
for tracking evaluation: nine and eight intervals of median \SI{18.8}{\second}
during which the ego is blind and a helper either covers the whole interval or
does not.

Because the occlusion and link-state columns are geometric predictions from a
map, each is checked against a measured signal that did not enter its
computation (Table~\ref{tab:validation}). Predicted occlusion is negatively
correlated with the evidence returned on both of \texttt{mobile\_1}'s
modalities; predicted link clearance is positively correlated with measured
RSSI on both mobile nodes, and the clearance-derived partition separates median
RSSI by \SI{8}{\decibel} and \SI{11}{\decibel}.

\begin{table}[t]
\centering
\caption{Validation of the geometric labels against independent measurement;
$\rho$ is the Spearman rank correlation over in-FoV observations.
$^\dagger$ Uninformative rather than negative: \texttt{mobile\_2} exceeded
$\tau$ in 2\,\% of its observations, leaving no variation to correlate against;
its evidence chain was validated along range instead (\num{861} depth points in
the box at \SI{2.4}{\metre}, none beyond \SI{7}{\metre}).}
\label{tab:validation}
\small
\begin{tabular}{@{}lll@{}}
\toprule
Label & Independent check & Result \\
\midrule
Occlusion (\texttt{m1}) & Ouster / ZED points in box & $\rho = -0.46$ / $-0.44$ \\
Occlusion (\texttt{m2}) & RealSense depth points in box & $\rho = -0.13$\,$^\dagger$ \\
Link clearance (\texttt{m1}) & measured RSSI & $\rho = +0.73$; LOS $-35$ vs NLOS $-43$\,dBm \\
Link clearance (\texttt{m2}) & measured RSSI & $\rho = +0.67$; LOS $-39$ vs NLOS $-50$\,dBm \\
\bottomrule
\end{tabular}
\end{table}
